\paragraph{The classical height theorem and the cusp limit.}
For a projective integral variety $X$ of dimension $n$ over a number
field $K$, normalize arithmetic intersections by $[K:\mathbb Q]$ and
point heights by the absolute degree of their field. Write
\[
 D(\bar L)=\deg_LX,\qquad
 I(\bar L)=\frac{\widehat{\deg}(\bar L^{n+1})}{[K:\mathbb Q]}.
\]
Zhang's successive-minima theorem, in its relative-semipositive form
\cite[Theorems A.4.2 and A.5.2]{yuan-zhang-adelic}, implies
\begin{equation}\label{geo-classical-height-lower}
 h_{\bar L}(P)\ge a\quad(P\in X(\overline K))
 \quad\Longrightarrow\quad
 I(\bar L)\ge(n+1)aD(\bar L).
\end{equation}
Here $L$ is ample. For a smooth Hermitian model, relative semipositivity
means semipositive curvature and nonnegative degrees on curves in
finite fibres; the metric is conjugation invariant. Nonnegative
horizontal point heights are not assumed. The point-height lower
bound makes every successive minimum at least $a$. Rational tensor
powers are included by homogeneity.

\begin{proposition}[Reciprocal perturbations suffice]
\label{geo-reciprocal-height-limit}
Fix a component and its prescribed cutoff and metrics, and finite
real numbers $a,\gamma,u_B$.
Let $\bar L_{m,j}$, $m,j\ge1$, be prescribed smooth Hermitian
$\mathbb Q$-line bundles on projective models of $X$, with perturbation
$1/m$ and smoothing index $j$. Suppose that, for each $m$ and all
sufficiently large $j$, they satisfy the positivity and ampleness
conditions above and
\[
             h_{\bar L_{m,j}}(P)\ge a-\gamma/m
                         \quad(P\in X(\overline K)).
\]
The threshold may depend on $m$ but must work for every point.
Suppose finite real $d_m,A_m,V_m$ satisfy
\[
 D(\bar L_{m,j})\longrightarrow d_m,\qquad
 \limsup_{j\to\infty}I(\bar L_{m,j})\le A_m+(n+1)V_m,
\]
and, for the original prescribed intersection $I_B$ and degree $d_B>0$,
\[
 A_m\longrightarrow I_B,\qquad d_m\longrightarrow d_B,\qquad
 V_m\longrightarrow u_Bd_B/2.
\]
Then
\[
                   \frac{I_B}{(n+1)d_B}\ge a-u_B/2.
\]
\end{proposition}

\begin{proof}
Apply \eqref{geo-classical-height-lower} and take the lower limit in
$j$. It is at most the assumed upper limit, so
\[
 (n+1)(a-\gamma/m)d_m\le A_m+(n+1)V_m.
\]
Let $m\to\infty$ and divide by $(n+1)d_B>0$. The argument permits
negative $a$, $u_B$ or $V_m$.
\end{proof}

For a finite nonempty set of actual components, suppose the same-tower test
identification is
\[
 I^{\rm test}_{B,Z}
 =s_{B,Z}(d_0+1)v_ZJ,\qquad
 s_{B,Z}=\frac{I_{B,Z}}{(n_Z+1)d_{B,Z}},\qquad
 v_Z,J>0,\quad d_0=fpq.
\]
With $v=\sum_Zv_Z$, the normalized total test is
\[
 \beta_B=\frac{\sum_ZI^{\rm test}_{B,Z}}{vfq(d_0+1)J}
         =\frac1{fq}\sum_Z\frac{v_Z}{v}s_{B,Z}.
\]
If Proposition~\ref{geo-reciprocal-height-limit} applies with
$a_Z=-C_hfq\Lambda_B$ and $u_{B,Z}\le C_ufq\Lambda_B$, then
\[
                   \beta_B\ge-(C_h+C_u/2)\Lambda_B.
\]
The signed component means may be retained before estimating them.
This is the input of Corollary~\ref{geo-tower-block-bound}. Its
remaining geometric premises are the prescribed positive
regularizations, the all-point height floor including the boundary,
and the one-sided intersection comparison and signed cusp limit.
A lower bound for the essential minimum alone does not imply
\eqref{geo-classical-height-lower}. Each datum is fixed before the
limits, so no convergence rate uniform in the growing family is
required for this deduction.
